\section*{Глава 1. Анализ предметной области}
\addcontentsline{toc}{section}{Глава 1. Анализ предметной области}
\stepcounter{section}

\subsection{Обзор существующих веб-сервисов для управления гардеробом}

На современном рынке цифровых решений для управления личным гардеробом выделяются приложения, предназначенные для каталогизации одежды, формирования образов и персонализированных рекомендаций пользователю. Такие сервисы позволяют систематизировать элементы гардероба и упрощают процесс подбора одежды в повседневной жизни.

Сравнительный анализ популярных решений представлен в таблице~\ref{tab:wardrobe_services}.
\begin{table}[ht!]
	\centering
	\footnotesize
	\setlength{\tabcolsep}{2pt}
	\caption{Сравнение существующих сервисов управления гардеробом}
	\label{tab:wardrobe_services}
	\begin{tabular}{|p{0.15\textwidth}|p{0.15\textwidth}|p{0.13\textwidth}|p{0.11\textwidth}|p{0.13\textwidth}|p{0.12\textwidth}|p{0.15\textwidth}|}
		\hline
		\textbf{Сервис} &
		\textbf{Каталогизация} &
		\textbf{Автоподбор} &
		\textbf{Веб-доступ} &
		\textbf{Изображения} &
		\textbf{Бесплатно} &
		\textbf{Расширяемость} \\ \hline
		
		Stylebook &
		Да &
		Нет &
		Нет &
		Да &
		Нет &
		Низкая \\ \hline
		
		Smart Closet &
		Да &
		Частично &
		Нет &
		Да &
		Частично &
		Средняя \\ \hline
		
		Cladwell &
		Да &
		Да &
		Нет &
		Да &
		Частично &
		Низкая \\ \hline
		
		Wardrobe AI &
		Да &
		Да &
		Да &
		Да &
		Частично &
		Средняя \\ \hline
		
		Pureple &
		Да &
		Частично &
		Нет &
		Да &
		Да &
		Низкая \\ \hline
		
	\end{tabular}
\end{table}

Проведённый анализ показал, что большинство существующих решений ориентировано преимущественно на мобильные платформы и имеет ограниченную поддержку веб-интерфейса. 

Сервисы Stylebook и Pureple предоставляют удобные инструменты каталогизации одежды, однако не поддерживают полноценный автоматический подбор образов. Cladwell и Wardrobe AI используют алгоритмы рекомендаций, но часть функциональности доступна только по подписке. 

Кроме того, существующие решения имеют ограниченные возможности расширения и адаптации под индивидуальные требования пользователя. Это подтверждает актуальность разработки собственного веб-приложения, ориентированного на браузерный доступ, хранение изображений одежды, формирование пользовательских образов и возможность дальнейшего расширения функциональности системы.

\subsection{Анализ программных инструментов разработки}

Выбор инструментов разработки определяется функциональными требованиями к веб-приложению для управления личным гардеробом. Система должна обеспечивать добавление, редактирование и удаление элементов одежды, хранение изображений, категоризацию предметов гардероба, формирование пользовательских образов и отображение данных через удобный веб-интерфейс.

\subsubsection*{\hspace{1.25cm}1.2.1 Выбор языка программирования}

Для разработки серверной части веб-приложения были рассмотрены несколько языков программирования: Python, JavaScript, PHP, Java, Go и C\#.

Python широко используется в веб-разработке и обладает простым синтаксисом, что ускоряет процесс создания приложения. Для Python существует большое количество веб-фреймворков, таких как Flask и Django, а также библиотек для работы с базами данных, изображениями и алгоритмами обработки данных. Это делает язык удобным для разработки приложения управления гардеробом и его дальнейшего расширения.

JavaScript может использоваться не только на стороне клиента, но и на сервере с помощью платформы Node.js. Такой подход позволяет разрабатывать клиентскую и серверную части на одном языке. Однако для реализации серверной логики и работы с базой данных требуется использование дополнительных инструментов и более сложная организация проекта.

PHP традиционно применяется для разработки веб-сайтов и веб-приложений. Он хорошо поддерживается большинством хостингов и имеет развитую экосистему. Однако в сравнении с Python он менее удобен для реализации алгоритмов анализа данных и дальнейшего расширения приложения в сторону интеллектуального подбора образов.

Java является надёжным языком для создания крупных корпоративных веб-систем. Она обеспечивает высокую производительность и строгую типизацию, однако требует более сложной структуры проекта и большего объёма кода, что делает её менее удобной для учебного проекта небольшого или среднего размера.

Go отличается высокой производительностью, простотой развертывания и хорошей поддержкой параллельной обработки. Он подходит для высоконагруженных серверных приложений, однако для рассматриваемого проекта его возможности являются избыточными, а экосистема для быстрой разработки пользовательских веб-приложений менее удобна по сравнению с Python.

C\# используется для разработки веб-приложений на платформе ASP.NET. Он обладает хорошей производительностью, строгой типизацией и развитой инфраструктурой. Однако использование C\# требует более сложной настройки среды разработки и лучше подходит для крупных проектов с корпоративной архитектурой.

Сравнение показало, что для разрабатываемого веб-приложения наиболее подходящим языком является \textbf{Python}. Он обеспечивает простоту разработки, хорошую поддержку веб-фреймворков, удобную работу с базами данных и возможность дальнейшего расширения функциональности, например добавления алгоритмов рекомендаций и анализа пользовательских данных.

\subsubsection*{\hspace{1.25cm}1.2.2 Выбор серверного фреймворка}

В качестве серверного фреймворка рассматривались Django, Flask и FastAPI.

Django предоставляет большое количество встроенных компонентов, включая систему аутентификации, административную панель и ORM. Однако данный фреймворк ориентирован на крупные проекты и обладает более сложной структурой.

FastAPI обеспечивает высокую производительность и удобную разработку API-приложений, однако в большей степени ориентирован на создание REST-сервисов и микросервисной архитектуры.

Flask является лёгким и гибким фреймворком, который предоставляет только основные инструменты для разработки веб-приложений. Это позволяет самостоятельно определить структуру проекта и реализовать только необходимую функциональность.

Для разрабатываемого приложения Flask является наиболее подходящим решением, поскольку он прост в освоении, подходит для учебного проекта и позволяет реализовать веб-приложение без избыточной архитектурной сложности.

Для разработки серверной части выбран фреймворк \textbf{Flask}. Он является лёгким и гибким инструментом, который подходит для веб-приложений небольшого и среднего размера. Flask позволяет реализовать маршруты для основных действий пользователя: добавления одежды, просмотра гардероба, редактирования записей, удаления элементов и создания образов. Также Flask удобно использовать в учебном проекте, так как он не навязывает сложную структуру и позволяет самостоятельно определить архитектуру приложения.

\subsubsection*{\hspace{1.25cm}1.2.3 Выбор базы данных}

При выборе системы управления базами данных рассматривались SQLite, MySQL и PostgreSQL.

SQLite является лёгкой встроенной базой данных и не требует отдельного сервера, однако её возможности ограничены при работе с многопользовательскими веб-приложениями.

MySQL широко используется в веб-разработке и обеспечивает высокую производительность, однако PostgreSQL предоставляет более развитые средства работы со сложными структурами данных и поддерживает расширенные возможности SQL.

PostgreSQL обеспечивает надёжное хранение данных, поддержку связей между таблицами и хорошо подходит для хранения информации о пользователях, элементах гардероба, категориях и сформированных образах.

Таким образом, \textbf{PostgreSQL} была выбрана как наиболее подходящая система управления базами данных для разрабатываемого веб-приложения.

\subsubsection*{\hspace{1.25cm}1.2.4 Выбор технологий интерфейса}

Для формирования пользовательского интерфейса используются \textbf{HTML} и \textbf{CSS}. Данные технологии подходят для отображения карточек одежды, форм добавления новых элементов, страниц со списком гардероба и готовых образов. CSS позволяет оформить интерфейс таким образом, чтобы пользователь мог удобно просматривать изображения одежды и взаимодействовать с элементами приложения.

Шаблонизатор \textbf{Jinja2} используется совместно с Flask для динамического формирования HTML-страниц. Он позволяет передавать данные из серверной части в интерфейс, например список предметов гардероба, категории одежды или сформированные комплекты. Это особенно важно для разрабатываемого приложения, поскольку содержимое страниц зависит от данных конкретного пользователя.

Таким образом, выбранный стек технологий соответствует задачам разрабатываемого веб-приложения.

\subsection{Формулировка цели и задач работы}

На основе анализа существующих решений и выявленных ограничений можно сформулировать цель курсовой работы.

\textbf{Цель работы} --- разработка веб-приложения для управления личным гардеробом, позволяющего пользователю добавлять, редактировать элементы одежды, систематизировать предметы гардероба по категориям и формировать пользовательские образы, а также выгружать готовые образы в виде изображений. Реализация приложения выполняется с использованием языка Python, фреймворка Flask, шаблонизатора Jinja2, технологий HTML и CSS, а также базы данных PostgreSQL.

Для достижения поставленной цели необходимо решить следующие задачи:

\begin{enumerate}
	\item Провести анализ существующих сервисов для управления личным гардеробом и формирования образов;
	\item Определить функциональные требования к разрабатываемому веб-приложению;
	\item Разработать архитектуру веб-приложения и структуру базы данных для хранения пользователей, элементов одежды, категорий и сформированных образов;
	\item Реализовать серверную часть приложения на языке Python с использованием фреймворка Flask;
	\item Разработать HTML-шаблоны пользовательского интерфейса с использованием Jinja2;
	\item Выполнить оформление страниц приложения с использованием CSS;
	\item Реализовать функциональность управления элементами гардероба: добавление, просмотр, редактирование и удаление записей;
	\item Реализовать загрузку и хранение изображений предметов одежды;
	\item Реализовать механизм формирования пользовательских образов на основе добавленных элементов гардероба;
	\item Провести тестирование разработанного веб-приложения.
    
\end{enumerate}
